\part{The findings in Grothendieck's language}
\section{An exact dictionary, not an attribution of new theorems}\label{sec:groth}
The new deductions above are restated here in the requested frameworks. A framework locator says where the language and operations live; it does not claim that a theorem about GCT occurs in Grothendieck's manuscripts. Where no archival proof or identification was verified, we say so explicitly.

\subsection{Schemes and group schemes over the integers}
The classical objects are group schemes $\GL_{r,\Z}$ and their polynomial representations on finite locally free modules. Tensor products and integral Schur constructions give maps such as
\[
 \GL(V)\times\GL(W)\longrightarrow\GL(V\otimes W),\qquad
 \GL(V)\longrightarrow\GL(S_\mu V).
\]
Representations of an affine group scheme are comodules of its \emph{commutative} coordinate Hopf algebra (SGA~3, I, \S4.2 and 4.7.2). Representations of diagonalizable groups correspond to graded modules (I, 4.7.3). These descriptions carry \eqref{eq:plethysm}, \fid{F11} and \fid{F13} into group-scheme language. Cote~90, scan~17, gives an authentic manuscript anchor for a flat affine group over a base, a group action, and equivariant quasi-coherent modules.

Over characteristic zero the representation categories of the reductive groups here are semisimple, so their actual classes form the positive monoid in the irreducible basis of $K_0$. This is the precise content of scalar nonnegativity. Over $\Z$ or in characteristic $p$, do not replace a filtration by a direct sum of simple modules. A divided Young symmetrizer is only an idempotent when its denominators are invertible; integral Schur functors must be constructed integrally. For the integral Specht-model representations used in the Kronecker and plethysm instances of \fid{F12}, tensor the free integral lattices with $\Z[1/|H|]$. The averaging idempotent then has a finite locally free direct-summand image, an effective $K_0$ class whose complex rank is the stated multiplicity. More generally this construction works over any coefficient ring carrying the representation and in which $|H|$ is invertible; an arbitrary complex representation need not descend to $\Z[1/|H|]$.

The nonstandard bialgebra preserves some comodule language but is not automatically a commutative affine group scheme. \fid{F11}'s matrix calculation proves why a reduction modulo $p$ supplies no automatic absolute Frobenius. This is a limitation of that construction, not a statement against separately defined quantum Frobenius maps.

\subsection{Representation rings, $\lambda$-rings and the $\gamma$-filtration}
For locally free sheaves, $\lambda^i[E]=[\wedge^iE]$ and Schur operations are the associated symmetric-function operations. On virtual classes the identity $\lambda_t(x-y)=\lambda_t(x)/\lambda_t(y)$ belongs to the group-completed ring. SGA~6, Expos\'e~V, \S\S2--3, contains the $\lambda$ and $\gamma$ formalism; cote~133, scan~18, explicitly displays symmetric/exterior generating functions. Cote~97, scan~4, concerns Chern classes of representations. These are the appropriate anchors for the formal calculation in \fid{F11}.

The representation ring example is an exact dictionary, not an analogy: $z$ is the standard one-dimensional comodule of $\mathbb G_m$; $z-1$ is a rank-zero virtual class; its exterior-power series gives $\lambda^2(z-1)=1-z$. The augmentation ideal and its $\gamma$-filtration retain characteristic-class information but introduce no order turning this virtual class into an actual representation. In particular the equality of two expressions in a $\lambda$-ring cannot replace the effective-simple-basis hypothesis of \fid{F07}.

For singular schemes, one must distinguish $K_0$ of vector bundles or perfect complexes from $G_0$ of coherent sheaves. The GCT graded pieces themselves are finite-dimensional representations, so their representation-$K_0$ classes are well-defined without identifying $K_0(X)$ and $G_0(X)$ on a singular orbit closure. Coherent Euler pushforward may involve alternating higher cohomology; an actual section module is more data than a virtual Euler class.

\subsection{Weil realizations and virtual versus effective traces}
In the trace formalism, a bounded complex of finite-dimensional Weil representations yields
\[
 Z(K,T)=\prod_i\det(1-TF\mid H^i(K))^{(-1)^{i+1}},\qquad
 \log Z(K,T)=\sum_{n\ge1}\Tr(F^n\mid K)\frac{T^n}{n}.
\]
The trace is an alternating trace. \fid{F02} bounds the surviving finite spectrum after equal roots have been grouped. \fid{F05}(b) says that $K_0$ and $Z(K,T)$ can erase identical spectra in adjacent degrees. An effective parity statement prevents this erasure; that is the exact role of (S2).

The geometric dictionary of \S\ref{sec:tracegeom} is the proof-carrying transfer of \fid{F01}--\fid{F04} and \fid{F17}. Character orthogonality is expressed by Artin--Schreier sheaves; the external tensor product realizes powers; compactly supported direct image realizes the sum; one-degree support fixes the alternating sign; Tate twist supplies the factor $q^n$; duality supplies the reciprocal eigenvalues. Framework locators are SGA~4, XVII, \S3 (compact supports), XVIII, 3.2.5 (smooth duality); SGA~5, III, 4.7 (Lefschetz--Verdier); and SGA~$4\frac12$, \emph{Rapport}, 3.2 (the trace formula). The exact ordinary weight estimates are Weil~II, 3.3.1 and 3.3.4--3.3.6. These general theorems are cited background; this run does not re-prove their machinery.

For a scheme over $\F_q$ and a closed point $x$ with residue field $\F_{q^{\deg x}}$, a Weil sheaf is pure of weight $w$ if every eigenvalue $\alpha$ at $x$ has
\[
 |\sigma(\alpha)|=q^{(\deg x)w/2}
\]
for every relevant complex embedding. For complexes the weight conditions are degree-shifted and involve duality; purity of a complex does not assert that every singular stalk has exact weight equal to its degree. The KL theorem uses the stronger stalkwise even/Tate statement \eqref{eq:KL}. This distinction translates \fid{F06} and prevents a false purity-to-coefficient inference.

\subsection{Orbit closures, Hilbert/Quot, and motives}
\fid{F15}'s kernel construction is the affine scheme-theoretic-image construction of EGA~I (1960), \S9.5. Its explicit $f=p$ example is already a group-scheme example over $\Spec\Z$. Projective closures with fixed Hilbert polynomial and coherent quotients belong to FGA, Bourbaki Expos\'e~221, \S3, Theorems~3.1--3.2, pp.~259--260. Fixing the Hilbert polynomial uses the chosen polarization, and a Quot functor specifies the coherent sheaf being quotiented. They do not turn coherent section multiplicities into intersection-cohomology multiplicities by representability alone.

For smooth projective schemes with algebraic correspondences, pure motives and their realizations provide another language for duality and projectors. The Tate realization in \fid{F16} is legitimate. By contrast a general singular orbit-closure intersection complex is not automatically a Chow motive furnished by SGA~4/5/6, and a quantum coordinate algebra acquires no such motive from its name. Cote~16, scans~35, 37 and~39, supplies a historical anchor for Lefschetz operators and primitive projectors; scans~47--49 discuss conditional correspondence/motive consequences, with conjectural letters interpreted in the letter's own notation. Their algebraic-projector and index assumptions cannot be replaced by the Euler sign $(-1)^i$.

\subsection{The constructive models and their exact comparison limits}
\paragraph{\fid{F18}: weight lattices and rational semigroups.}
Regard $g(\nu,(d,d),(d,d))$ as the coefficient of
$\det^d\boxtimes\det^d$ in the restriction of $S_\nu(V\otimes W)$,
where $V,W$ have rank two. This is an equality in the representation
$K_0$ of the group scheme $\GL(V)\times\GL(W)$ after passage to
characteristic zero. Characters of a diagonalizable group correspond
to gradings (SGA~3, I, 4.7.3). Integral difference-constraint arrays
would supply a particular saturated graded semigroup. The flooring
and dilation proof of \fid{F18} transfers unchanged to that weight
lattice. Its parity example requires rational-lattice or congruence
data instead. SGA~6's formal $\lambda$ identities do not supply those
data; cote~133, scans~18--21, is a vocabulary anchor rather than a
positive-rule theorem.

\paragraph{\fid{F19}, \fid{F21}, \fid{F25}: coherent algebra and its actual module.}
Here there is a specified scheme over a localization of
$\Spec\Z[v,v^{-1}]$: the finite flat center $S$ of \fid{F19}, with
the explicitly trivial bundle $E$ of ranks $1,1,2$.
The noncommutative algebra is $\Gamma(S,\End E)$, and the actual
right tensor module is $\Gamma(S,\mathcal Hom(E,\mathcal F))$ for
the displayed coherent branch sheaf $\mathcal F$. Thus the module
dictionary is an explicit action-preserving isomorphism, not just
an equality of characters. The matrices in \fid{F21} carry the proof
over to this sheaf language. Equivariant quasi-coherent modules are
the framework of cote~90, scan~17, and SGA~3, I, 4.7.2.

The branch modules $L_i$ are coherent but are not asserted to be
locally free or perfect over the singular center. Their natural
Grothendieck group is $G_0(C)$; one must not silently apply vector-bundle
$\lambda$ operations from $K_0(C)$ to them. The periodic resolution
in \fid{F25} computes their actual coherent derived Hom. Its
$R/(\Delta)$ records an order-two deformation invisible in the
constant coefficient realization. The regular algebra basis in that
finding is unrelated to any assertion that this singular coherent
category is semisimple. Its scalar $\lambda$ has compact-support motive
\[
 \mathbb L^{-2}(\mathbb L-1)^2(1+\mathbb L+\mathbb L^2)
\]
in the virtual Tate ring. The trace equality follows from
$\mathbb G_m^2\times\PP^2$ and the trace formula; its mixed signs are
cohomological Euler signs, not dimensions of one pure degree.

\paragraph{\fid{F22}: the \'etale topos and the structure sheaf retain different data.}
The finite maps $\Spec C_n\to\Spec C_m$ are universal
homeomorphisms. EGA~IV, 18.1.2, in the precise modern formulation
Stacks 04DZ/03SI, transports the entire ordinary coefficient-sheaf
topos over the trait. The proof in \fid{F22} explicitly identifies
the constant nearby-cycle triangle with
$\Q_\ell\to\Q_\ell^2\to\Q_\ell$; all three objects have scalar
weight-zero Frobenius. In contrast the conductor and crossed coherent
$\operatorname{Ext}^1$ retain $m$. This is an exact instance of the
difference between a coefficient topos and a ringed coherent space.
It does not forbid adding a separately specified nilpotent action to
a Weil sheaf. No manuscript statement is claimed to prove this
nonstandard deformation comparison.

\paragraph{\fid{F20}: effective Schur operations and Tate coefficient motives.}
On a split affine Grassmannian, the primitive Young projector is an
actual idempotent in a rational symmetric tensor category. Its image
is an effective object, and its class is the specified Schur operation
on the locally defined representation. The modern motivic Satake
theorems cited in \fid{F20} identify its coefficient motive as
$\mathbf1(-h)^{\oplus a}$, not merely an alternating virtual class.
The ordinary coefficient realization has scalar Frobenius $q^h$.
This supplies an actual motive in these classical cases; it does not
assert the existence of a Chow motive for every singular orbit closure.
The source of the motivic theorem is Richarz--Scholbach, not SGA~6
or cote~16. The latter's scans~35,37,39 supply the historical language
of algebraic projectors. Canonical bases and efficient rank labels
remain separate information even when the motive is Tate.

\paragraph{\fid{F23}: induction in the effective representation monoid.}
A signed orbit is the representation of a constant finite group
scheme induced from a one-dimensional stabilizer representation.
Over $\Z[1/d!]$ with the relevant integral lattices, finite-group
averaging splits invariant projectors; over characteristic zero
the representation category is semisimple. The direct-sum equality
of \fid{F23} is effective in representation $K_0$, and its Young
stabilizer specialization is the ordinary symmetric/exterior Pieri
operation. The proof by orbit modules and restriction of Schur
functors is unchanged in this language. The pure Tate $\GL_2$
example shows why the $S_d$ correspondence action is additional to
the weight and motive. Cotes~90/17 and~133/19--21 give the respective
module and finite-set-operation anchors; neither provides the missing
universal signed basis.

\paragraph{\fid{F24}: a signed class can still retain a largest eigenvalue.}
For a bounded Weil complex, group equal roots in its alternating trace.
Its visible multiplicities are integers, possibly of both signs.
The external tensor product maps this class to its power in the
group ring of eigenvalues. The augmentation-and-Vandermonde argument
in \fid{F24} proves that an even power retains a positive top-radius
coefficient. This transfers the proof exactly to the triangulated
$K_0$ trace formalism of SGA~4/5, without requiring one input degree.
The negative visible path-output spectrum and any reciprocal visible
duality are still hypotheses. A zero virtual class may erase actual
cohomology, and finite extensions can erase a signed top spectrum;
the two countermodels specify those limits.

\paragraph{\fid{F26}--\fid{F27}: localization, bounded complexes, and finite traces.}
Finite Weil cohomology gives rational resolvents and zeta factors
$\det(1-TF)$, including their alternating products. This is the
finite determinant mechanism underlying the trace formalism; the
non-Laurent rational grading function in \fid{F26} cannot be realized
by such a fixed finite trace in all extensions. Its recurrence proof
is independent of any particular choice of category. The generic
Fredholm determinant formulas visible in cote~1, scans~77--78,
do not identify it with a finite Weil object or a $\Xi$-function.

Bounded complexes in the Hecke category do realize the signed Euler
class $P$. Their four basis coordinates prove the one-parity
obstruction in \fid{F27}. The rank-one bimodule comes from the
finite flat correspondence
$\Spec A\times_{\Spec A^s}\Spec A$. Separately, the coefficient
Hecke algebra $H=R[c]/(c^2-\delta c)$ normalizes to $R^2$
(and its tensor square to $R^4$); this separates its character
points and yields the other positive cone.
The actual four-point Tate construction retains the rank-one action
after the stated localization. Algebraic localization, coherent
normalization, and a trace-compatible tensor realization are three
specified operations here; none is inferred from the others.


\subsection{Finding-by-finding transfer table}
{\small\begin{longtable}{p{.12\linewidth}p{.48\linewidth}p{.32\linewidth}}
\toprule IDs&Exact Grothendieck-framework form&Locator/anchor and limit\\\midrule\endhead
F01--F04, F17&Trace of external tensor powers on the path hypersurface; parity of $R\Gamma_c$; Tate duality; optional finite base extension for each $k$. The calculations of \S\ref{sec:main} apply after the trace-formula identification.&SGA~4 XVII \S3, XVIII 3.2.5; SGA~5 III 4.7; Weil~II 3.3.1. Cote~16 provides only the distinct projector/duality framework, not this theorem.\\
F05&The zero virtual class $[\Q_\ell(-1)]-[\Q_\ell(-1)]$ has invisible large eigenvalues. A pure arithmetic Weil module can be a nontrivial extension.&Triangulated $K_0$; BBD 5.3.8 and its following example, pp.~138--139. No archival proof of the new examples is asserted.\\
F06--F07&Effective graded simple classes yield nonnegative multiplicities; a positive phase and Lefschetz maps are additional. Euler classes alone do not suffice.&BBD 5.3.2, 5.4.5, 5.4.10; cote~16, scans~35, 37, 39 for operator language.\\
F08&The rank vector bundles on two points form a special $\lambda$-ring; an effective integral basis need not be its simple positive basis.&SGA~6 V \S\S2--3; cote~133, scan~18. The entire counterexample is proved here.\\
F09--F10&An algebra inclusion need not preserve the effective monoid of a split $K_0$ categorification; exclusive monomial pivots give a sufficient alternative.&Split Grothendieck groups, effective objects, grading shifts; cote~90, scan~17 for equivariant-module context. We do not attribute the nonstandard inclusion claim to the manuscripts.\\
F11&Virtual exterior powers fail effectivity; noncommutative $p$th powers fail additivity, so do not define a scheme Frobenius.&SGA~3 I 4.2, 4.7.2--3; SGA~6 V \S\S2--3; cotes~97/4 and~133/18.\\
F12--F13&For the integral Specht models, invariant projectors are direct summands over $\Z[1/|H|]$; positive restriction matrices compose in representation $K_0$.&SGA~3 I 4.7.2--3; cote~90/17. Actual theta/GCT functor is still absent.\\
F14&Frobenius-stable subquotients inherit a weight; endomorphism spectra are ratios.&Weil representations; BBD 5.3.1 for perverse subquotient weight bounds. Realization/stability is still required.\\
F15&Kernel-defined orbit schemes over $\Z$; coherent polarization and linearization retained in $K_0/G_0$; finite-field base change is separate.&EGA~I \S9.5; FGA~221, 3.1--3.2; cote~90/17 for group actions.\\
F16&The centered polynomial of $\mathbf1\oplus\mathbb L^{\oplus2}\oplus\mathbb L^2$ realizes the rank-one product coefficient; its action comparison fails.&Projective bundle/K\"unneth calculations; cote~16 correspondence language. Pure Tate pieces, not an ungraded pure weight-zero module.\\
F18&Homogeneous weight-lattice difference arrays force ordinary support saturation; a rational half-point retains parity.&SGA~3 I4.7.3 for gradings; exact proof F18; VI Definition1.2.4 is a distinct residuewise notion.\\
F19, F21, F25&A finite flat center, a vector bundle endomorphism algebra and the actual coherent right tensor module; coherent Ext measures the collision.&Cote90/17; affine scheme/module dictionary; explicit complete proofs in these findings. Branch modules belong to coherent $G_0$.\\
F20&Primitive Schur projectors on convolution have effective pure Tate coefficient motives with fixed twists.&Richarz--Scholbach 4.1,4.8,5.9,6.2--6.5; BBD weight formalism; cote16 projector vocabulary only.\\
F22&Universal homeomorphisms preserve the ordinary coefficient topos while conductor and coherent Ext vary.&EGA~IV18.1.2; Stacks04DZ,03SI,0A4E.\\
F23, F28&Effective induced representation classes yield orbit/Pieri formulas when their stabilizer hypotheses hold; an isolated rectangular hook excludes that Young-orbit model.&SGA~3 representations; cote133/19--21 only a finite-set-operation anchor. IP Proposition6.4 supplies the hook input.\\
F24&Even external tensor powers of visible signed Euler spectra retain a positive largest-radius coefficient.&SGA~4/5 trace formalism; full proof F24, with powered-spectrum visibility and visible duality.\\
F26--F27&Finite Weil resolvents are rational; non-Laurent extension evaluations are excluded. Bounded Hecke complexes and normalized character schemes have different effective cones.&Split/triangulated $K_0$; finite normalization; cote1/77--78 gives generic Fredholm vocabulary, not this finite trace identity.\\\bottomrule
\end{longtable}}

\paragraph{\fid{F28}: rectangular projectors in representation $K_0$.}
The unprojected multiplicity module is the actual constant-group
representation $[(a^b)]\otimes[(a^b)]$. A proposed monomial orbit
basis would be an effective sum of induced one-dimensional Young
subgroup representations. Exterior powers of the standard
representation compute their consecutive hook supports. The
isolated hook supplied by IP Proposition~6.4 contradicts such an
effective sum; this proof carries verbatim to representation $K_0$.
A virtual sum may cancel those hooks and is a different assertion.
The known Manivel algebra instead has a torus-graded monomial basis,
with the displayed determinant-character comparison. This belongs
to the affine-group and weight-grading framework of SGA~3, I,
4.7.2--3 and cote~90, scan~17. Cote~133's operations on finite sets
do not identify these two different positive structures.


\paragraph{\fid{F29}: flag schemes and orbit dimensions.}
The action is that of the split group scheme $SL_n\times SL_n$ on
the partial flag scheme of its tensor product representation,
base-changed to $\C$. These flag schemes are iterated relative
Grassmannians, hence instances of the locally free Quot framework
in FGA, Expos\'e~221, \S3, Theorems~3.1--3.2. Their filtered
frame bundles give the Borel torsors used in the proof; the group
actions belong to SGA~3, I, \S4. The complete dimension and
stabilizer calculation proves only the absence of an open orbit for
this action. Cote~90, scan~17, anchors equivariant affine-module
language, not an attribution of this orbit theorem to the archive.

\paragraph{\fid{F30}: vector bundles and coherent commutants under specialization.}
In the language of schemes, $M_R$ is a vector bundle on $\Spec R$
with an action of a finite coherent algebra. Images of the regular
idempotents above are coherent direct summands and hence vector
subbundles. The regular mixing is an automorphism of the underlying
bundle; it does not commute with the whole algebra action. The coherent
commutant is the kernel of the finitely many linear commutation
equations for the generators and is the scalar sheaf in this case.
Its formation does not commute with the special fiber: that fiber
has a second classical projector. This is an explicit base-change
phenomenon for coherent endomorphisms, compatible with Grothendieck's
scheme and module language. It supplies neither a Frobenius action nor
a weight statement, and purity alone does not decide which algebra
equivariance a proposed comparison should retain.

The framework is the affine module/sheaf equivalence of EGA~I,
\S9.5, and equivariant modules in SGA~3, I, 4.7.2--3; cote~90,
scan~17, is the audited archive anchor. Every needed idempotent
image, scalar commutant and special-fiber calculation is proved
in \fid{F30}; no general base-change theorem is asserted without
its hypotheses.

\paragraph{\fid{F31} in Grothendieck's language.}
The integral Grassmannian is the projective Quot scheme of
rank-$d$ locally free quotients of $\mathcal O_{\Z}^{n+d}$,
with its universal sequence and Pl\"ucker line. FGA, Expos\'e~221,
\S3, Theorems~3.1--3.2, pp.~259--260, supplies the representability
framework. Its split $\GL_{n+d}$ action and the cocharacter used
above belong to SGA~3, I, 4.7.2--3. The cokernel is an actual
Weil object and the projector an actual rational correspondence,
with effective class $g_r[\mathbb L^r]$. The inverse exists here
because of the full matrix Hom spaces of the Tate motive, while
the integral Schubert cone need not be preserved. Cote~16,
scans~35, 37 and~39, gives the audited Lefschetz/primitive-projector
vocabulary; these leaves are not asserted to contain this
Kronecker identification or its proof.


\paragraph{\fid{F32}: completed Tate classes and the limits of visible duality.}
For a countable direct sum of Tate pieces whose Euler trace has a
convergent nonnegative absolute moment, the theorem specifies exactly
the discrete completion being used. It is not an element of ordinary
finite $K_0$ merely by notation. Tensor multiplication is legitimate
because the absolute-moment estimate proves convergence and finite
grouping at each nonzero root. The formal binomial expansion of
$(1+\mathbb L^{-1})^{-4}$ in the actual trace has such a realization
at the level of convergent Euler spectra; its inverse Tate degrees
are unbounded. Dualization sends them to unbounded positive degrees,
outside this bounded-spectrum class. Thus an effective or signed
completed class still needs a specified duality and trace functor.
The finite trace framework is SGA~4\,$\frac12$, \emph{Rapport},
3.2, and SGA~5, III, 4.7. Cote~1, scans~77--78, records the
audited generic Fredholm trace/determinant vocabulary; neither it
nor finite $K_0$ supplies this countable extension. Its convergence
and reciprocal-finiteness proof is given in \fid{F32} itself.

\paragraph{\fid{F35} in Grothendieck's language.}
The ordinary split group schemes and quotient presentation are
explicit. Positive-dimensional stabilizers are retained, so this
construction uses the Artin-stack extension of the six operations,
not the bounded constructible pushforward theorem for schemes.
Its rational mass and infinite alternating Tate trace are connected
by the precisely cited Laszlo--Olsson/Sun formalism. The effective
cohomology groups in \eqref{eq:stackcompactgroups} have positive
dimensions, while their grouped Euler coefficients have both signs.
The duality exists but does not preserve the chosen convergent
geometric-Frobenius trace class. The audited manuscript determinant
and projector vocabulary is only a framework anchor; no specific
stack, nonstandard action, or proof here is attributed to those leaves.

\paragraph{\fid{F38} in Grothendieck's language.}
The projective-line decomposition and the stack's actual Tate
cohomology distinguish a positive object from its virtual Euler
class. A rational mass can invert an Euler trace while leaving
unbounded nonzero cohomology. Thus localization of a coefficient
ring or of a convergent trace algebra does not itself provide an
invertible geometric object. The split $\GL_2$ group scheme,
projective bundle formula, and modern stack six operations are
the precise frameworks used here. No archive leaf is claimed to
contain this construction or to identify it with a nonstandard
canonical basis. A different categorical localization requires a
specified new realization and its own comparison theorem.

\paragraph{\fid{F39} in Grothendieck's language.}
The singular piece is the actual affine scheme
$\Spec k[\epsilon]/\epsilon^2$, and matrix Morita equivalence
expresses its coherent category without changing the defect.
Smoothness here asks for a perfect diagonal bimodule, a perfect-complex
condition distinct from finite-dimensionality or ordinary $\ell$-adic
purity. The displayed resolution proves its failure directly. EGA~I's
affine scheme language and the perfect-complex framework of SGA~6
provide the classical dictionary; the smooth proper dg and
noncommutative-motive extension is explicitly the modern cited theory.
The audited cote~16, scans~48--49, concerns conditional semisimplicity
and Weil statements in its own motivic framework; it is not a source
for this dg construction. A new realization and basis comparison is
required before its Frobenius can address the GCT coefficients.


\paragraph{\fid{F34}: effective representation classes and growth of graded algebras.}
The actual highest-weight multiplicities are ranks of effective Schur
summands, in the characteristic-zero representation $K_0$ of $\GL_2$.
Their exact positive simplex has dimension $k-2$; purity cannot alter
that growth degree. The polynomial bound fails only for the specified
compact encoding. Locally free Schur operations and their virtual
extensions use the SGA~3 representation and SGA~6 $\lambda$-ring
frameworks. Cote~133, scans~18--21, anchors those operations, not the
computational bound. The dimension and output-length arguments are
proved explicitly in \fid{F34}; they are not archival theorems.

\paragraph{\fid{F37}: effective multiplicity objects with compressed labels.}
The known rectangular identity identifies the marked-tree count with
the rank of the primitive Tate summand constructed in \fid{F31}.
Compression is a bijection on the finite indexing sets for each input,
so it changes neither that effective representation class nor its rank.
It gives a polynomial witness encoding under the listed-row convention;
it does not identify the trees with a canonical cycle basis or make their
recursion into a theta correspondence. Cote~133, scans~19--21, is a
framework anchor for finite-set operations, and cote~16, scans~35,37,39,
for primitive projectors. The exact counting input is PPS, Theorem~4.5;
the new encoding argument is the proof of \fid{F37} itself.



\paragraph{\fid{F36}: coherent Frobenius, normalization and descent.}

This is a calculation of finite maps, coherent modules, relative
Frobenius and derived tensor products on an explicit singular scheme.
The relative map from collision order $mp$ to order $m$ is a finite
universal homeomorphism, like those in \fid{F22}; its coherent cokernel
is nevertheless nonzero. Ordinary constant \'etale traces lose this
length. The normalization has unit coherent Frobenius and preserves
the original algebra action after restriction, while pushforward
through the collision does not preserve the unit property.


The affine module/sheaf language is EGA~I, \S9.5; the action and
equivariant-module language is SGA~3, I, 4.7.2--3, with the audited
cote~90, scan~17, as its archival anchor. The derived calculations
are the explicit resolutions in the proof, not an invocation of
flat base change at the singular point. No archival identification
of coherent absolute Frobenius with a nonstandard Weil realization
is asserted.


\paragraph{\fid{F34}, constructive supplement: the actual graded ray scheme.}
The actual scheme $\Spec B$ is affine $(k-1)$-space with positive
grading $2,\ldots,k$; $\operatorname{Proj}B$ is the corresponding
weighted projective scheme. The degree-$N$ graded coordinate piece
is the highest-weight multiplicity space just proved, so its Hilbert
series is $\prod_{i=2}^k(1-z^i)^{-1}$. The exponent simplex has
dimension $k-2$. These are exact statements about the affine
coordinate ring and its grading, without an identification of
coherent Hilbert multiplicities with arbitrary constructible
cohomology. The ordinary affine and graded-scheme dictionary applies;
no arithmetic weight, designated nonstandard basis, or archival
claim is supplied by this coordinate computation alone.

The affine dictionary is EGA~I, \S9.5, and the highest-weight representation dictionary uses SGA~3, I, \S4. The separate source anchors for Schur operations are recorded above; no complexity bound is ascribed to those texts.


\paragraph{\fid{F40} in Grothendieck's language.}
Affine endomorphism algebras, idempotent corners and perfect
diagonals give the precise coherent dictionary. The explicit
finite diagonal resolution proves smoothness of a larger category,
while its projective objects retain the singular corner and its
actual module. This uses the perfect-complex language associated
with SGA~6 and the explicitly cited modern smooth dg and
noncommutative-motive theory. The exceptional decomposition
explains the unit Frobenius spectrum without turning the corner
subcategory into an admissible motivic summand. No unpublished
archive leaf is claimed to contain this construction. A relative
order, a realization of the designated morphisms and a comparison
of the required bases remain separate mathematical tasks.

Exact classical framework locators are EGA~I, \S9.5, and SGA~6, Expos\'e~V, \S\S2--3, the latter retained as background with the source-access qualification in the archival audit. Cote~16, scans~48--49, is a vocabulary anchor for conditional motivic semisimplicity and Weil statements, not for this smooth dg construction.


\paragraph{\fid{F33} in Grothendieck's language.}
The geometric object is the ordinary split group scheme $SL_m$.
Its rational Tate-motive decomposition supplies actual Weil
summands and scheme Frobenius, so the equality of dimensions
has a precise realization. The relevant frameworks are SGA~3's
split group schemes and the explicitly cited modern motive
theorem. The generator degree $2i-1$ and weight $2i$ remain
different gradings. No manuscript leaf is claimed to contain
this hook formula or a nonstandard basis comparison.


The group-scheme language is SGA~3, I, \S4; cote~90, scan~17, is the audited group-action anchor. The motive decomposition itself is Merkurjev, Corollary~5.5, and not an assertion from the archive.


\paragraph{\fid{F41} in Grothendieck's language.}
The effective Tate summands and their weights are those of the
actual stack calculation in \fid{F35}. Eigenspace restriction
belongs to ordinary Weil linear algebra, and a finite weight
window restores the finiteness premise of \fid{F26}.
The chosen fixed-parameter action demonstrates that a geometric
coefficient object and an extension-wise parameter comparison
are distinct requirements. No action on the whole stack, or
archive attribution for such an action, is asserted.

The finite trace framework is SGA~4\,$\frac12$, \emph{Rapport}, 3.2, and SGA~5, III, 4.7. The infinite stack input uses the precisely cited modern extension; no bounded scheme theorem is applied to the unbounded object.


\paragraph{\fid{F42} in Grothendieck's language.}
The finite-resolvent obstruction concerns one finite realization
and its fixed Frobenius, whereas the Tate coefficient objects
use tensor powers with growing spectra. The fixed-parameter
construction is an exact action in the category of point Weil
objects, with named generators and tensor compatibility; it is
not an extension-wise quantum grading comparison. Scalar
division, an idempotent image, and inversion of a tensor object
are distinct operations here. The remaining path comparison
must preserve the actual operations and control its boundary
complex; no such comparison is attributed to an archive leaf.

The exact trace-framework locators are SGA~4\,$\frac12$, \emph{Rapport}, 3.2, and SGA~5, III, 4.7; the tensor/action language uses SGA~3, I, 4.7.2--3. Cote~90, scan~17, supplies the audited action vocabulary, not the comparison axiom.

\section{Archival audit and the Fable leads}\label{sec:archive}
The original Montpellier PDFs add a cover. For the documents inspected here, a bottom-stamped archive scan $s$ is one-based PDF page $s+1$. Catalogue titles and scanned contents are distinguished. The following checks were visual readings of the specified leaves, not certifications of full manuscript proofs.

{\small\begin{longtable}{p{.19\linewidth}p{.36\linewidth}p{.37\linewidth}}
\toprule Cote and scan&Verified content&Consequence for this investigation\\\midrule\endhead
90, scan 17 (PDF 18)&Flat affine group over a base; group action; quasi-coherent equivariant modules.&Authentic anchor for the representation/group-scheme dictionary.\\
97, scan 4 (PDF 5)&Chern classes of regular representations and products over characters.&Characteristic-class anchor, not a plethysm-positivity theorem.\\
133, scan 18 (PDF 19)&Symmetric and exterior generating operations, including $\sigma_t(E)\lambda_{-t}(E)=1$.&Direct language anchor for the $\lambda$-ring calculation.\\
133, scans 19--21 (PDF 20--22)&Virtual finite $G$-sets and operations; scan 20 explicitly names plethysm.&These finite-set operations must not be identified with a proved representation-theoretic positive rule.\\
16, scans 35, 37, 39 (PDF 36, 38, 40)&Handwritten Lefschetz/primitive-summand/projector formulas, including hyperplane-section comparisons.&Relevant standard-conjecture framework. No identification with Alp\"oge's path sign is proved.\\
16, scans 34, 36, 38, 40&Reused typescript backs; scan 34 starts Karoubi's \emph{Cohomologie des cat\'egories de Banach}.&The Fable range ``34--40'' must be refined; these backs are not Grothendieck's handwritten Hodge-index proof.\\
16, scans 47--49 (PDF 48--50)&Letter to Coates: correspondences/duality; on scan 48, $C(X)$ plus $I(X\times X)$ implies correspondence-ring semisimplicity; on scan 49, deduction of Weil from $C$ and $A$, in the letter's notation.&Conditional standard-conjecture statements. The semisimplicity passage is scan 48, not scan 50.\\
16, scan 50 (PDF 51)&Separate handwritten endomorphism-algebra lemma.&Not the location of the preceding semisimplicity sentence.\\
1, archival scans 77--78 (PDF78--79), handwritten folios79--80&Exterior-power/trace-norm estimates; scan78 has a Fredholm determinant product and trace sum.&Generic operator-determinant framework is visible; the heavily revised manuscript proof is not certified.\\
1, archival scans109--110 (PDF110--111), handwritten folios111--112&Entire functions, order/genus and canonical products.&No $\Xi$ identity or RH-equivalent positive operator was verified.\\
1, literal archival scans79--80 and111--112 (PDF80--81,112--113)&Corrected determinant calculations, convex/supporting functions, and entire-function or positive-sequence convergence material.&Both numbering interpretations were checked; these are not interchangeable leaves.\\\bottomrule
\end{longtable}}
Cote~133's catalogue title concerns Witt groups, quadratic forms and formal groups; ``plethysm'' is a content lead, not its official title. Cote~60 was found in the catalogue and inspected at thumbnail level, but no precise theorem from it is used.

The Fable suggestions are treated as follows. Suggestion~3's norm estimate is proved and extended in \fid{F01}; its literal positive-Fourier-operator identification fails by \fid{F05}(e). Suggestion~4's identification of the cohomological sign with a polarized evanescent intersection form remains unproved: parity in a trace formula and a positive-definite form are different structures. Suggestion~8 is made precise by \fid{F13}, conditional on an actual functor and endpoint identification. No assertion about the unprovided earlier runs or their theorem IDs is imported.

Some published-source checks have a narrower scope than the archive checks. The original SGA~6 scan could not be visually retrieved here; Expos\'e~V, \S\S2--3, is carried as background. The specific locator V, Lemma~2.2.1, p.~315, is independently cited by Altman--Kleiman \cite[\S2.10]{AK}; it is a locator check, not a re-proof of SGA~6. The modern transcription consulted as a finding aid is not used as technical authority. All elementary $\lambda$ identities used in the paper have their displayed direct proofs.

\part{Unproved bridges, attempted routes, and current state}
\section{What a sufficient nonstandard package would have to contain}\label{sec:gates}
There are no unchecked affirmative ``new theorems'' in this part. The following is a sufficient research specification, not a quotation from Mulmuley and not a proved construction.

\begin{enumerate}
\item \gate{R: realization and normalization.} For each relevant nonstandard coefficient construct a finite-field Weil object or a defined categorical replacement, with explicit Frobenius, trace formula, degrees and weights. Identify $v$ with the correct grading, and specify signs, half-powers and rational rescalings.
\item \gate{S: spectral sign and duality.} Verify the pointwise input bound, the finite-spectrum identification, the output sign for unbounded $k$, and controlled background terms. F24 allows signed visible inputs at even $k$; its finite-extension version requires retained top-radius visibility. Supply either the Tate self-duality of \fid{F03} or the dual system of \fid{F17}. These inputs imply the stated purity theorem by the checked proof.
\item \gate{E: effective basis and operations.} Identify the normalized coefficient with a graded multiplicity of actual simple objects and prove compatibility with multiplication/comultiplication. Purity alone does not give this. The fixed inherited ambient cone is ruled out by \fid{F09}; F27 restores the actual rank-one character action in a different localized cone. Neither identifies VIII's basis.
\item \gate{L: Lefschetz and defects.} Prove the appropriate Lefschetz maps for unimodality. In general plethysm, define the nonflatness correction and prove the quantitative error bounds of the intended version of VIII 2.9/3.6. A spectral error bound such as \fid{F04} is not automatically that coefficient-vector error bound.
\item \gate{C: canonical and comparison data.} Prove basis existence/globalization, cell and quasi-cell claims, and the required specialization. For general plethysm retain VII's duality/reciprocity hypotheses; for Kronecker do not re-open the already proved revised-IV duality.
\item \gate{P: effective labels under a stated encoding.} Specify the input convention before asking for polynomial witnesses. Finding~\fid{F37} settles the equal-rectangle, two-row-output witness problem with explicitly listed rows. Finding~\fid{F34} excludes a general \(\#P\) claim under fully compact binary plethysm encoding. The attached homogeneous PH1 is already refuted by the cited Kahle--Micha{\l}ek theorem; it is not an open requirement to be proved. Any replacement using inhomogeneous polytopes or another counting model needs its own exact formulation, size bounds and algorithm. SH or the appropriate index bound remains separate.
\item \gate{O: actual obstructions.} Establish a strict multiplicity comparison for the orbit closures in an appropriate asymptotic family. Positive formulas and efficient decisions do not themselves produce that inequality.
\end{enumerate}
The spectral theorems remove an estimate once the explicit trace and sign hypotheses are supplied. F32 extends the upper-radius argument to a specified countable completion, but bounded visible reciprocal duality then forces finiteness. F19--F21 and F27 construct specified coherent or ordinary motivic realizations. F36 constructs action-compatible coherent Frobenius models and computes their descent defect. F35 closes the scalar stack-realization question for the actual normalized word trace; its dual does not satisfy the needed second bound. F38 excludes the proposed object-level tensor inverse in the usual derived Weil realization, and F39 excludes the direct nonsmooth collision category from the smooth proper noncommutative Weil framework. F40 repairs that gate by an explicit smooth proper corner category, retaining the exact specialized action; its whole-category TP spectrum is pure but does not yet realize the designated coefficient morphisms. These results narrow particular realizations; the designated based operations and path-output comparison remain open. Replacing all seven by the single phrase ``nonstandard RH implies positivity'' hides independent claims.

For completeness, an effective-label claim has a simple formal implication: if the labels $\alpha$ in \eqref{eq:comparison} and witnesses for $m_\lambda^\alpha,n_\pi^\alpha$ have polynomial-length unique encodings and polynomial-time tests, then triples consisting of a label and the two witnesses count the sum of products. A nondeterministic machine guesses and checks the triple, giving a \(\#P\) function. This is just the definition of \(\#P\); neither the encodings nor a replacement for the refuted PH1 follows from the abstract multiplicity equality. No complexity-class collapse or lower bound is asserted.

\section{Approaches attempted and their exact outcomes}
{\small\begin{longtable}{p{.23\linewidth}p{.40\linewidth}p{.29\linewidth}}
\toprule Approach&Proved outcome or failure point&Next unresolved input\\\midrule\endhead
Literal source formulation&No finite-field nonstandard RH statement in the supplied versions; actual signed/near-positive targets located.&Choose a precise R/E target.\\
Path orthogonality and repeated Parseval&F01 gives the bound, but the constant trace example defeats purity without spectral signs.&S1/S2 realization.\\
Tensor powers and inverse roots&F02--F04 remove the fixed exponent loss; F17 allows separate finite extensions and a dual system.&Prove the concentration/error and duality hypotheses.\\
Fourier-positive operator&The bare Fourier-after-nonnegative-multiplication claim is false, F05(e).&A different operator and a proved identity to the coefficient.\\
Ordinary ambient purity&Stable subquotients inherit purity, F14. The actual inherited positive Hecke cone has only scalar image, F09.&A new effective realization/heart or different character map.\\
Change of rank-one positive model&F16 loses the actual action; F27 restores it in a localized four-character mixed Tate cone.&Identify the designated basis and operations; the repaired object is not pure of one weight.\\
Positive supports/crystals&Monomial pivots suffice, F10. Global positive disjoint supports and local disjointness alone fail.&Global exclusive pivots or another valid basis argument.\\
Orbit/Hilbert/Quot geometry&Integral models and ordinary IC purity exist; polarized coherent multiplicities are not determined by underlying Weil cohomology, F15.&A coherent-to-constructible multiplicity comparison.\\
Determinants and Gram forms&Positive projectors already express the ranks, F12; chosen-basis coefficient positivity fails in F08.&Image basis, independence and operation compatibility.\\
Theta and interlacing chains&F13 gives path/polytope counts for specified exact functors.&Identify the actual theta chain and endpoints with GCT restriction.\\
$\lambda$-rings and motives&Formal operations and Tate examples translate exactly; virtual effectivity and algebraic projectors remain extra.&Effective classes and the required realization, not only ring identities.\\\bottomrule
\end{longtable}}

\section{States of the five requested tasks}
\begin{enumerate}
\item \textbf{Statement needed.} The source-location and target audit is completed as A01 and \S\ref{sec:targets}. The demand for exact nonstandard Frobenius objects/weights cannot be met from those papers because they are not specified. The replacement package in \S\ref{sec:gates} is explicitly our sufficient specification.
\item \textbf{Template.} The standard geometric dictionary and every missing role are stated in \S\ref{sec:KL}. The actual rank-three coherent algebra and right tensor module are supplied by F19/F21; F20 supplies ordinary pure Tate classical coefficients. The designated nonstandard based realization remains open at R/E.
\item \textbf{Sign-only mechanism.} F01--F04, F17, F24 and the admissible countable extension F32 are proved and checked; omitted-hypothesis versions are refuted by the displayed countermodels. F26 excludes the finite normalized all-extension word dictionary. F35 supplies a genuine stack scalar realization with explicit cohomology, while F38 separates its reciprocal traces from actual tensor inverses. The specified main transfer remains open at R/S/E; a scalar mass identity is not a construction of the operators, path output, dual bound or designated basis.
\item \textbf{Positivity reached.} F31 gives actual primitive Tate motives for equal rectangles with two-row output. F37 gives polynomially verifiable compressed witnesses for the known marked-tree rule under the listed-row encoding. F34 proves the separate compact dimension and output-length obstructions. The supplied homogeneous PH1 is refuted by a cited prior theorem, including an equal-rectangle example. General designated canonical-basis positivity and strict orbit-closure multiplicity inequalities remain open at E/L/C/O; replacement effective rules require the revised P gate.
\item \textbf{Theta positivity.} F12--F13 and F23 state the exact projector/path consequences; F18 and F28 delimit integral-interlacing and Young-orbit versions. The theta/GCT functor and basis identification remain open. F08 refutes the generic determinant-to-basis-positivity shortcut.
\end{enumerate}

Research is continuing on the remaining gates under the scheduled-review rule. The present checkpoint is not a final return. No result here proves or refutes the intended full nonstandard positivity program, or $P\ne NP$.

\appendix
\section{Verification and reproducibility}
The independent spectral audit reconstructed F01, reproved F02 by Ces\`aro orthogonality rather than poles, and checked the two-limit argument of F03. The initial all-conjugate duality gap was recorded and repaired before the final theorem. Later full-paper reviews checked each proof, including the actual Hecke-cone example, the positive-pivot countermodels, the categorical and theta statements, and the scheme/base-change examples. Exact dates, reviewer runs, model identity, exposed effort information and check scopes are in \texttt{CHECK\_EVENTS.tsv}.

The file \texttt{check\_path\_bound.py} checks 54 small scalar/vector examples, including complex weights, for the exact path identity and the estimate. These floating-point checks are supplementary; the proof is F01. The optional exact Murnaghan--Nakayama script and character certificate compute a small rectangular coefficient and check character orthogonality; they are not used in any proof here and are not a Lean certificate. The source manifest records the unchanged attachments by SHA-256. Source reads and whole-proof checks are distinguished from cited background theorems and from formal verification.

The dated ledger records actual changes of gates. Scheduled checkpoints are due only at 45-minute intervals; scheduled progress reviews begin only after approximately three hours. The metadata records which such deadlines, if any, occurred; no review or checkpoint is invented to imply elapsed work that did not occur.

\begin{thebibliography}{99}
\bibitem{Alpoge} L.~Alp\"oge, \emph{Remark}, supplied four-page note, 20 November 2026; \url{https://alpo.ge/remark.pdf}. Main locators: p.~3 footnote~3 and steps~3--5; p.~4 steps~7--8. Supplied formalization entry points: \texttt{ALPOGE\_Main.lean}, \texttt{ALPOGE\_Statements.lean}; scan record \texttt{ALPOGE\_LEAN\_SCAN.tsv}. Results accepted as stipulated; missing imported development not executed.
\bibitem{IV} J.~Blasiak, K.~D.~Mulmuley and M.~Sohoni, \emph{Geometric Complexity Theory IV: nonstandard quantum group for the Kronecker problem}, arXiv:cs/0703110v4, 2013. \url{https://arxiv.org/abs/cs/0703110v4}. Exact statement locators are given in the text; attached source \texttt{Apr13KroneckerGCT4.tex}.
\bibitem{VI} K.~D.~Mulmuley, \emph{Geometric Complexity Theory VI: the flip via saturated and positive integer programming in representation theory and algebraic geometry}, arXiv:0704.0229v4, 2009. \url{https://arxiv.org/abs/0704.0229v4}.
\bibitem{VII} K.~D.~Mulmuley, \emph{Geometric Complexity Theory VII: Nonstandard quantum group for the plethysm problem}, arXiv:0709.0749v2, 2008. \url{https://arxiv.org/abs/0709.0749v2}.
\bibitem{VIII} K.~D.~Mulmuley, \emph{Geometric Complexity Theory VIII: On canonical bases for the nonstandard quantum groups}, arXiv:0709.0751v2, 2008. \url{https://arxiv.org/abs/0709.0751v2}.
\bibitem{IP} C.~Ikenmeyer and G.~Panova, \emph{Rectangular Kronecker coefficients and plethysms in geometric complexity theory}, arXiv:1512.03798v2, 2017; especially Theorems~1.4, 1.6, 1.8, 1.10. \url{https://arxiv.org/abs/1512.03798v2}.
\bibitem{BIP} P.~B\"urgisser, C.~Ikenmeyer and G.~Panova, \emph{No occurrence obstructions in geometric complexity theory}, arXiv:1604.06431v3, 2018; Theorem~1.4, \S1(c), Proposition~3.3, equation~(4.1), Proposition~4.5, Theorem~4.7. \url{https://arxiv.org/abs/1604.06431v3}.
\bibitem{KL} D.~Kazhdan and G.~Lusztig, \emph{Schubert varieties and Poincar\'e duality}, in \emph{Geometry of the Laplace Operator}, Proc. Sympos. Pure Math.~36, 1980, pp.~185--203; Theorem~4.3. The exact normalized form used here is checked through Haines's primary proof below.
\bibitem{Haines} T.~J.~Haines, \emph{A proof of the Kazhdan--Lusztig purity theorem via the decomposition theorem of BBD}; Theorems~2.0.1, 5.0.4, Proposition~3.0.2, \S4. \url{https://www.math.umd.edu/~tjh/KL_purity1.pdf}.
\bibitem{WeilII} P.~Deligne, \emph{La conjecture de Weil II}, Publ. Math. IH\'ES~52 (1980), pp.~137--252; Definitions~1.2.2 and~6.2.2--6.2.4, Theorem~3.3.1 and Corollaries~3.3.4--3.3.6. \url{https://www.numdam.org/item/PMIHES_1980__52__137_0/}.
\bibitem{BBD} A.~Beilinson, J.~Bernstein and P.~Deligne, \emph{Faisceaux pervers}, Ast\'erisque~100 (1982). Exact locators: 5.1.2, 5.1.8, 5.1.14, 5.3.1--5.3.2 (including the second proof), 5.3.8 and the following example, 5.4.4--5.4.5, Remark~5.4.9 (p.~143), 5.4.10 (p.~144), 6.2.5. \url{https://www.numdam.org/item/AST_1982__100__1_0/}.
\bibitem{Molev} A.~I.~Molev, \emph{Gelfand--Tsetlin bases for classical Lie algebras}, Theorem~2.1, equation~(2.2), patterns~(2.4); author-hosted preprint, printed pp.~7--9. \url{https://www.maths.usyd.edu.au/u/pubs/publist/preprints/2002/molev-18.pdf}.
\bibitem{Milne} J.~S.~Milne, \emph{Lectures on \'Etale Cohomology}, v2.21, 2013; Example~16.3, p.~109; Theorem~23.2, pp.~139--140. \url{https://www.jmilne.org/math/CourseNotes/LEC.pdf}.
\bibitem{EGA} A.~Grothendieck and J.~Dieudonn\'e, \emph{\'El\'ements de g\'eom\'etrie alg\'ebrique I: Le langage des sch\'emas}, Publ. Math. IH\'ES~4 (1960), \S9.5. Framework locator, with the affine instance proved in F15. \url{https://www.numdam.org/item/PMIHES_1960__4__5_0/}.
\bibitem{SGA3} M.~Demazure and A.~Grothendieck (eds.), \emph{Sch\'emas en groupes (SGA 3)}, Expos\'e~I by M.~Demazure, \S4.2, 4.7.2--4.7.3. Checked in the Gille--Polo re-edition, respecting explicitly marked editorial additions. \url{https://wstein.org/sga/SGA3/Exp1-20jan2010.pdf}.
\bibitem{SGA4} M.~Artin, A.~Grothendieck and J.-L.~Verdier (eds.), \emph{Th\'eorie des topos et cohomologie \'etale des sch\'emas (SGA 4)}, Expos\'es~XVII, \S3, and XVIII, 3.2.5. Background framework references.
\bibitem{SGA45} P.~Deligne, \emph{Cohomologie \'etale (SGA $4\frac12$)}, LNM~569, 1977; \emph{Rapport sur la formule des traces}, 3.2. Background trace-formula reference.
\bibitem{SGA5} A.~Grothendieck (ed.), \emph{Cohomologie \(\ell\)-adique et fonctions L (SGA 5)}, LNM~589, 1977; Expos\'e~III, 4.7. Background Lefschetz--Verdier reference.
\bibitem{SGA6} P.~Berthelot, A.~Grothendieck and L.~Illusie (eds.), \emph{Th\'eorie des intersections et th\'eor\`eme de Riemann--Roch (SGA 6)}, LNM~225, 1971; Expos\'e~V by P.~Berthelot, \S\S2--3. Original scan not visually verified in this run; used as background, not a claimed fresh source-proof check.
\bibitem{FGA} A.~Grothendieck, \emph{Techniques de construction et th\'eor\`emes d'existence en g\'eom\'etrie alg\'ebrique IV: les sch\'emas de Hilbert}, S\'eminaire Bourbaki, Expos\'e~221 (1961), \S3, Theorems~3.1--3.2, pp.~259--260. \url{https://www.numdam.org/article/SB_1960-1961__6__249_0.pdf}.
\bibitem{AK} A.~B.~Altman and S.~L.~Kleiman, \emph{Foundations of the theory of Fano schemes}, Compositio Math.~34 (1977), pp.~3--47; \S2.10, used only to corroborate the SGA~6 locator. \url{https://www.numdam.org/item/CM_1977__34_1_3_0/}.
\bibitem{Archive} Universit\'e de Montpellier, \emph{Archives Grothendieck}, cotes~1, 16, 90, 97, 133, scans specified in \S\ref{sec:archive}. Catalogue: \url{https://grothendieck.umontpellier.fr/archives-grothendieck/}; original files \url{https://grothendieck.umontpellier.fr/16.pdf}, \url{https://grothendieck.umontpellier.fr/90.pdf}, \url{https://grothendieck.umontpellier.fr/97.pdf}, \url{https://grothendieck.umontpellier.fr/133.pdf}. Cote~1 was checked at eight selected leaves, with both scan and handwritten-folio numbering as recorded in the text.
\bibitem{Fable} Claude Fable 5.2, effort max as reported in its supplied file, \emph{FABLE\_PROGRAM\_SUGGESTIONS\_20261120.md} and \emph{\ldots\_EXPLAINED.md}, suggestions~3, 4 and~8. Unchecked leads, not cited theorems.
\bibitem{Richarz} T.~Richarz, \emph{A new approach to the geometric Satake equivalence}, Documenta Mathematica 19 (2014), 209--246. Theorem~3.1, Proposition~4.1, Theorem~5.1, Lemma~5.4, and their proofs; the numbering differs from arXiv 1207.5314v2. \url{https://ems.press/content/serial-article-files/26224?nt=1}.
\bibitem{MV} I.~Mirkovi\'c and K.~Vilonen, \emph{Geometric Langlands duality and representations of algebraic groups over commutative rings}, Annals of Mathematics 166 (2007), 95--143; arXiv math/0401222v5 includes the later Appendix~B correction. Theorem~7.3, Lemma~7.1, Corollary~7.4. \url{https://arxiv.org/abs/math/0401222}.
\bibitem{StacksNil} The Stacks Project, Theorem~59.45.2, tag~04DZ, and Proposition~59.45.4, tag~03SI, topological invariance of the \'etale site/topos; Lemma~59.91.9, tag~0A4E, finite base change. \url{https://stacks.math.columbia.edu/tag/03SI}.
\bibitem{RS} T.~Richarz and J.~Scholbach, \emph{The motivic Satake equivalence}, Math. Ann.~380 (2021), 1595--1653; arXiv:1909.08322v3. Theorem~4.1, Propositions~4.8 and~5.9, Definition~6.2, Lemma~6.3, Corollary~6.4, Lemma~6.5, Theorem~6.8. \url{https://arxiv.org/abs/1909.08322}.
\bibitem{GW} R.~Goodman and N.~R.~Wallach, \emph{Symmetry, Representations, and Invariants}, author-hosted Chapter~4, Theorem~4.1.18 (Hermite reciprocity), p.~191. \url{https://mathweb.ucsd.edu/~nwallach/chapter4.pdf}.
\bibitem{DK} R.~W.~Donley, Jr., and W.~G.~Kim, \emph{A rational theory of Clebsch--Gordan coefficients}, arXiv:1707.03022v2; Proposition~3.1, equation~(3.4), p.~3; Proposition~8.1, equation~(8.2), p.~11. \url{https://arxiv.org/abs/1707.03022}.
\bibitem{Pieri} A.~Postnikov, \emph{MIT 18.315 lecture notes}, scribed by C.~Lian, 2014; Theorem~7.10, p.~16. Used for the classical horizontal and vertical Pieri rules. \url{https://web.mit.edu/clian/www/315_notes.pdf}.
\bibitem{EW} B.~Elias and G.~Williamson, \emph{The Hodge theory of Soergel bimodules}, Annals of Mathematics~180 (2014), 1089--1136; \S3.5, Theorem~3.5, pp.~1107--1108; Theorem~1.1, p.~1093. \url{https://annals.math.princeton.edu/wp-content/uploads/annals-v180-n3-p06-p.pdf}.
\bibitem{Manivel} L.~Manivel, \emph{A note on certain Kronecker coefficients}, arXiv:0809.3710v1 (2008); Theorem~1, p.~1; \S3, pp.~3--4, for the four-generator invariant algebra. \url{https://arxiv.org/abs/0809.3710}.

\bibitem{BlasiakRep} J.~Blasiak, \emph{Representation theory of the nonstandard Hecke algebra}, arXiv:1201.2209v2 (2012); Theorem~5.1, p.~14; Proposition~5.2, pp.~14--15; Definition~6.1 and Remark~6.2, p.~23. PDF equation~(25) has a different number in the converted HTML. \url{https://arxiv.org/abs/1201.2209v2}.
\bibitem{PPBounds} I.~Pak and G.~Panova, \emph{Bounds on certain classes of Kronecker and $q$-binomial coefficients}, author manuscript dated 13 October 2016; Lemma~1.3, p.~2. \url{https://www.math.ucla.edu/~pak/papers/QBC23.pdf}.
\bibitem{Brosnan} P.~Brosnan, \emph{On motivic decompositions arising from the method of Bia\l ynicki--Birula}, Invent. Math.~161 (2005), 91--111; inspected author PDF \S2, pp.~3--4, Theorems~3.3 and~3.5, pp.~5--6. \url{https://math.umd.edu/~pbrosnan/Papers/bb.pdf}.


\bibitem{KM} T.~Kahle and M.~Micha{\l}ek, \emph{Plethysm and lattice point counting}, Electronic Journal of Combinatorics~25(1) (2018), P1.41; Theorem~5 and Remark~6, p.~4; formulation history, pp.~2--4; rectangular identification, p.~6. \url{https://www.combinatorics.org/ojs/index.php/eljc/article/download/v25i1p41/pdf/}.
\bibitem{PPS} I.~Pak, G.~Panova and J.~P.~Swanson, \emph{A combinatorial interpretation for certain plethysm and Kronecker coefficients}, arXiv:2511.02312v3, uploaded 27 July 2026, title-page date 28 July 2026; Lemma~1.3, p.~2; \S2.2, p.~5; Definition~3.2 and Proposition~3.4, p.~6; Lemma~4.2, p.~9; Definition~4.4 and Theorem~4.5, p.~10. \url{https://arxiv.org/abs/2511.02312v3}.
\bibitem{SunStacks} S.~Sun, \emph{L-series of Artin stacks over finite fields}, Algebra \& Number Theory~6 (2012), 47--122; Theorem~1.1, p.~49; Definition~4.1, p.~84; Theorem~4.2, p.~85; Theorem~4.8 and \S4.8.2, pp.~89--90; Example~7.4, p.~109. \url{https://msp.org/ant/2012/6-1/ant-v6-n1-p03-s.pdf}.
\bibitem{LO2} Y.~Laszlo and M.~Olsson, \emph{The six operations for sheaves on Artin stacks II: Adic coefficients}, Publ. Math. IH\'ES~107 (2008), 169--210; \S9.1, Lemmas~9.1.1--9.1.2, p.~198; Theorem~11.0.14, p.~204. \url{https://numdam.org/item/10.1007/s10240-008-0012-5.pdf}.
\bibitem{TabuadaWeil} G.~Tabuada, \emph{Noncommutative Weil conjecture}, inspected version arXiv:1808.00950v7 (2019); \S1, pp.~2--3, Theorem~1.5; \S5.1, pp.~13--14; \S6.1, pp.~15--16; \S7, equations~(7.1)--(7.4), p.~17; Example~12.3, p.~24; Theorem~2.12, p.~8, and proof equation~(12.6), p.~26; \S6.2, equation~(6.9), p.~16. \url{https://arxiv.org/abs/1808.00950v7}.

\bibitem{MerkurjevSL} A.~S.~Merkurjev, \emph{Motivic decomposition of certain special linear groups}, inspected author PDF, \S5, Corollary~5.3 and notation, p.~7; Corollary~5.5, p.~8. \url{https://www.math.uni-bielefeld.de/LAG/man/530.pdf}.
\bibitem{Run} This run, \texttt{GCT\_NSRH\_20261121}; findings A01 and F01--F42, with agent source memoranda, proof reviews, computations and \texttt{CHECK\_EVENTS.tsv} in the final package. This is the only current-run identifier used.
\end{thebibliography}
\end{document}
